\documentclass[10pt,a4paper]{amsart}
\usepackage[margin=19mm]{geometry}
\usepackage{amsmath,amssymb,mathtools}
\usepackage[scr=rsfs,cal=euler]{mathalfa}
\usepackage{xfrac}
\usepackage[unicode,hidelinks,bookmarks=false]{hyperref}
\newcommand{\ie}{i.e.\ }
\newcommand{\cf}{cf.\ }
\newcommand{\resp}{respectively\ }
\newcommand{\Subsection}{\subsection}
\providecommand{\nobreakdash}{\nobreak\hbox{-}}
\setlength{\emergencystretch}{2em}
\multlinegap0pt
\usepackage{bm}
\usepackage{bbm}
\usepackage{dsfont}
\usepackage{xspace}
\usepackage{cancel}
\usepackage{mathtools}
\usepackage[shortlabels]{enumitem}
\usepackage{amsthm}
\makeatletter
\@ifundefined{lema}{\newtheorem{lema}{Lemma}}{}
\@ifundefined{propo}{\newtheorem{propo}{Proposition}}{}
\makeatother
\DeclareMathOperator{\Diff}{Diff}
\DeclareMathOperator{\Hol}{Hol}
\DeclareMathOperator{\Sing}{Sing}
\newcommand{\C}{\mathbb C}
\newcommand{\F}{\mathcal F}
\newcommand{\N}{\mathbb N}
\newcommand{\OX}{\mathcal O}
\title[Holomorphic and Formal First Integrals for Foliations of Cod]{Holomorphic and Formal First Integrals for Foliations of Codimension One on Complex Analytic Space Germs}
\author{Victor Le\'on, Bruno Sc\'ardua}
\date{Source: arXiv:2608.06491v1 (CC BY 4.0)}
\begin{document}
\maketitle

\noindent\emph{Source and licence.} Holomorphic and Formal First Integrals for Foliations of Codimension One on Complex Analytic Space Germs. Version arXiv:2608.06491v1 (CC BY 4.0), distributed under the Creative Commons Attribution 4.0 International licence, as recorded in the arXiv licence metadata for this exact version and shown on the abstract page https://arxiv.org/abs/2608.06491v1. Full rights notice: licenses/arxiv-2608.06491-CC-BY-4.0.txt.\par\smallskip

\noindent\textit{Editorial scope.} Counted unit: the Propo whose source label is \texttt{Proposition:6} in the article (source0.tex lines 2265--2294, source label \texttt{Proposition:6}), delivered together with the article's own complete original proof of it (source0.tex lines 2296--2414, one \texttt{proof} environment of 119 lines). No mathematics is written, repaired or re-derived: statement and proof are the article's own text line for line. Required context retained uncounted: lemma:6 (lines 2150--2160); lemma:7 (lines 2237--2239). Every definition, display and label of those lines is retained unchanged. Omitted from this package and not redistributed: 1--2149, 2161--2236, 2240--2264, 2295--2295, 2415--2668. Nothing inside a retained excerpt is omitted or altered: every retained body is delivered exactly as the article's lines are written, so no omission receipt of any kind accompanies this package. Citations: the retained excerpts carry no citation command at all, so no bibliography entry of the article is transcribed and no reference is redistributed. Reference adaptation: none was needed and none was made; every reference inside the delivered unit resolves inside it, so no unresolved reference or citation is produced. Printed slips kept literal: no printed word, symbol, hypothesis or number of the counted statement or of its proof was altered. Layout and class: the article's own class is \texttt{amsart} with packages including lmodern, amsmath, amssymb, amsfonts, amsthm, mathtools, mathrsfs, geometry, xcolor, cite, hyperref, tikz, tikz-cd, enumitem; the wrapper is the lane's pinned \texttt{amsart} editorial wrapper at 10pt on A4, which restates the theorem-like environments the retained text needs and adds the article's own macro definitions that the retained text needs (\texttt{\textbackslash C}, \texttt{\textbackslash Diff}, \texttt{\textbackslash F}, \texttt{\textbackslash Hol}, \texttt{\textbackslash N}, \texttt{\textbackslash OX}, \texttt{\textbackslash Sing}), with extra wrapper packages (bm, bbm, dsfont, xspace, cancel, mathtools, enumitem). The delivered numbering is the unit's own. Assets: no retained span contains a figure environment, a graphics-inclusion command, a PDF-page inclusion, a table or a TikZ picture; no image file is copied into this package, the package declares no asset dependency and no image file is redistributed with it. Licence: version arXiv:2608.06491v1 of this article is distributed under the Creative Commons Attribution 4.0 International licence, as recorded in the arXiv licence metadata for this exact version and shown on the abstract page https://arxiv.org/abs/2608.06491v1. A full rights notice is delivered at licenses/arxiv-2608.06491-CC-BY-4.0.txt. Characters: every retained excerpt is pure ASCII, so the delivered engine, XeLaTeX with its default Latin Modern text font, reports no missing character.
\begin{lema}
\label{lemma:6}\normalfont
Let $G<\Diff(\C,0)$ be a group of germs of one-variable complex analytic
\textup{(not necessarily finitely generated)} diffeomorphisms at $0\in\C$.
If there is a non-constant formal function
$\hat f\in \widehat{\OX}_{\C,0}$, $\hat f(0)=0$, such that
\[
        \hat f\circ \varphi = \hat f, \qquad \forall\varphi\in G,
\]
then $G$ is finite.
\end{lema}

\begin{lema}
\label{lemma:7}\normalfont A holonomy map $h\in \Hol(\widetilde\F,L_j,\Sigma_j,q_j)$ leaves invariant the restriction $\widehat F|_{\Sigma_j}$.
\end{lema}

\begin{propo}\label{Proposition:6}\normalfont
Let $\F$ be a germ of holomorphic foliation on a germ of complex analytic
normal space $(X,0)$ of dimension two. Assume that $\F$ admits a formal
first integral $\hat f\in\widehat\OX_{X,0}$. Then:
\begin{enumerate}[label=\textup{(\arabic*)}]
\item $\F$ is non-dicritical;
\item $\widetilde\F$ exhibits no saddle-nodes;
\item every singularity $\tilde p\in\Sing(\widetilde\F)\subset E$ is
analytically linearizable of the form
\[
       kx\,dy+l y\,dx=0,\qquad k,l\in\N,\quad \gcd(k,l)=1,
\]
where
\[
       \{y=0\}\subset E\subset \{xy=0\};
\]
\item writing $E=\displaystyle\bigcup_{j=1}^r E_j$ in compact components,
the holonomy group of every component is finite:
\[
H(E_j)\subset\Diff(\Sigma_j,q_j),
\]
where
\[
       q_j\in E_j\setminus\Sing(\widetilde\F), \qquad
       \Sigma_j\pitchfork E, \qquad
       \Sigma_j\cap E=\Sigma_j\cap E_j=\{q_j\},
\]
for every choice of $q_j,\Sigma_j$ and every $j\in\{1,\ldots,r\}$.
\end{enumerate}
\end{propo}

\begin{proof}
\textup{(1)} After subtracting its constant term, we may assume that
\[
\widehat f(0)=0.
\]
Put
\[
\widehat F:=\pi^*\widehat f.
\]
Then $\widehat F$ vanishes along every component of
$E=\pi^{-1}(0)$.

Fix a component $E_j$ and choose a generic smooth point
$q\in E_j$. Let $x=0$ be a reduced local equation for $E_j$ near
$q$. Since $\widehat F$ is nonzero, we may write
\[
\widehat F=x^m\widehat U,
\qquad m\geq1,
\]
where $\widehat U$ is not divisible by $x$. Let $\widetilde v$ be a
local generator of $\widetilde\F$. Since $\widehat F$ is a formal
first integral,
\[
0=\widehat{\widetilde v}(\widehat F)
=x^{m-1}\left(
m\widehat U\,\widetilde v(x)
+x\widehat{\widetilde v}(\widehat U)
\right).
\]
Reducing the resulting identity modulo $x$, we obtain
\[
m\,(\widehat U|_{E_j})\,
(\widetilde v(x)|_{E_j})=0.
\]
At a generic point of $E_j$, $\widehat U|_{E_j}\neq0$. Hence
\[
\widetilde v(x)|_{E_j}=0,
\]
on a nonempty open subset of $E_j$. Since
$\widetilde v(x)|_{E_j}$ is holomorphic on the regular part of the
irreducible component $E_j$, the identity theorem gives
\[
\widetilde v(x)|_{E_j}\equiv0.
\]
Equivalently, $\widetilde v(x)\in(x)$ along $E_j$, and therefore
$\widetilde v$ is tangent to the entire component $E_j$. Thus every
component of $E$ is $\widetilde\F$-invariant, and consequently
$\F$ is non-dicritical.

\textup{(2)} Given a singularity
$\tilde p\in\Sing(\widetilde\F)\subset E$, the formal function
$(\pi^*\hat f)_{\tilde p}\in\widehat\OX_{\widetilde X,\tilde p}$ is a
formal first integral for the germ $\widetilde\F_{\tilde p}$ of the
foliation $\widetilde\F$ at $\tilde p\in\widetilde X$.
By Mattei--Moussu, Theorem A, page 472, the germ $\widetilde\F_{\tilde p}$
admits a nonconstant holomorphic first integral. Since $\pi$ is a
resolution of the pair $((X,0),\F)$, every singularity of
$\widetilde\F$ is simple, and hence is either a reduced nondegenerate
singularity or a saddle-node. A saddle-node cannot admit a nonconstant
holomorphic first integral. Therefore $\widetilde\F$ has no saddle-nodes,
which proves \textup{(2)}.

It remains to identify the nondegenerate singularities. A reduced
nondegenerate singularity admitting a holomorphic first integral has two
analytic separatrices and a negative rational quotient of eigenvalues.
Thus, after choosing relatively prime integers $k,l\in\N$ and suitable
analytic coordinates, a primitive holomorphic first integral can be
written as
\[
F(x,y)=x^l y^k.
\]
Consequently, the foliation is analytically equivalent to
\[
kx\,dy+l y\,dx=0.
\]
Since the exceptional divisor is invariant and has normal crossings, the
coordinates may be chosen so that one of its local branches is
$\{y=0\}$ and
\[
\{y=0\}\subset E\subset\{xy=0\}.
\]
This proves \textup{(3)}.

We now prove \textup{(4)}. Following the notation introduced above in the
proof, we consider a component $E_j$ of $E$ and choose a generic
nonsingular point $q_j\in E_j\setminus\Sing(\widetilde\F)$,
and a transverse disk $\Sigma_j$ centered at
$\{q_j\}=\Sigma_j\cap E_j$.
Put $\widehat F:=\pi^*\widehat f$ and denote by
$\widehat F|_{\Sigma_j}$ its restriction to the formal transverse
section $(\Sigma_j,q_j)$.
For every $h\in H(E_j)$, Lemma~\ref{lemma:7} gives
\[
       \bigl(\widehat F|_{\Sigma_j}\bigr)\circ h
       =\widehat F|_{\Sigma_j}.
\]
Choose a local coordinate $x$ transverse to $E_j$, so that
$E_j=\{x=0\}$, and write
\[
\widehat F=x^{m_j}\widehat U,
\qquad m_j\geq1,
\]
where $\widehat U$ is not divisible by $x$. We choose $q_j$ outside
the zero set of the leading coefficient $\widehat U|_{E_j}$. Hence
$\widehat U|_{\Sigma_j}$ is a unit, and
\[
\operatorname{ord}_{q_j}
\bigl(\widehat F|_{\Sigma_j}\bigr)=m_j\geq1.
\]
Thus $\widehat F|_{\Sigma_j}$ is nonzero, nonconstant, and vanishes at
$q_j$. Lemma~\ref{lemma:6} implies that
$H(E_j)$ is a finite group. For any
other regular point $q_j'\in E_j\setminus\Sing(\widetilde\F)$, transport
along the connected leaf
$E_j\setminus\Sing(\widetilde\F)$ conjugates the corresponding
holonomy group to the one just considered. Therefore the
holonomy group is finite for every choice of a regular
base point and transverse section on $E_j$.
\end{proof}

\end{document}
